\documentclass[11pt,a4paper]{article}

\usepackage[T1]{fontenc}
\usepackage[utf8]{inputenc}
\usepackage{lmodern}
\usepackage{amssymb}
\usepackage{microtype}
\usepackage[a4paper, top=2.4cm, bottom=2.4cm, left=2.7cm, right=2.7cm]{geometry}
\usepackage{xcolor}
\usepackage{booktabs}
\usepackage{array}
\usepackage{enumitem}
\usepackage{titlesec}
\usepackage{parskip}
\usepackage{fancyhdr}
\usepackage{tcolorbox}
\usepackage{hyperref}
\tcbuselibrary{skins,breakable}

\definecolor{qblue}{HTML}{1a2a3a}
\definecolor{qaccent}{HTML}{3d5a80}
\definecolor{qlight}{HTML}{eef3f7}
\definecolor{qwarn}{HTML}{a5350d}
\definecolor{qwarnbg}{HTML}{fbeee7}
\definecolor{qgrey}{HTML}{7a7a7a}
\definecolor{qcode}{HTML}{2d2d40}
\definecolor{qcodebg}{HTML}{f4f4f6}

\hypersetup{colorlinks=true, linkcolor=qaccent, urlcolor=qaccent,
  pdftitle={qmeas v2 User Manual}, pdfauthor={qmeas project}}

\titleformat{\section}{\normalfont\Large\bfseries\color{qblue}}{\thesection}{0.6em}{}
\titleformat{\subsection}{\normalfont\large\bfseries\color{qaccent}}{\thesubsection}{0.6em}{}
\titlespacing{\section}{0pt}{1.4em}{0.5em}

\pagestyle{fancy}
\fancyhf{}
\fancyhead[L]{\small\color{qgrey}qmeas v2}
\fancyhead[R]{\small\color{qgrey}\leftmark}
\fancyfoot[C]{\small\color{qgrey}\thepage}
\renewcommand{\headrulewidth}{0.4pt}

\newtcolorbox{tip}{breakable, skin=enhanced, colback=qlight, colframe=qaccent,
  boxrule=0.4pt, left=2mm, right=2mm, top=1mm, bottom=1mm,
  fonttitle=\bfseries, title=Tip}
\newtcolorbox{warn}{breakable, skin=enhanced, colback=qwarnbg, colframe=qwarn,
  boxrule=0.4pt, left=2mm, right=2mm, top=1mm, bottom=1mm,
  fonttitle=\bfseries, title=Watch out}
\newtcolorbox{note}{breakable, skin=enhanced, colback=qwarnbg, colframe=qwarn,
  boxrule=0.4pt, left=2mm, right=2mm, top=1mm, bottom=1mm,
  fonttitle=\bfseries, title=Note}

\newcommand{\code}[1]{{\ttfamily\color{qcode}#1}}

\newtcolorbox{gridex}[1]{skin=enhanced, colback=qcodebg, colframe=qgrey!50,
  boxrule=0.5pt, arc=0.6mm, left=2mm, right=2mm, top=1.4mm, bottom=1.4mm,
  fontupper=\ttfamily\small, title={#1}, fonttitle=\bfseries\small\color{qblue},
  coltitle=qblue, colbacktitle=qcodebg, titlerule=0pt}

\newtcolorbox{cmd}{skin=enhanced, colback=qcodebg, colframe=qgrey!45,
  boxrule=0.5pt, arc=0.6mm, left=2.4mm, right=2mm, top=1.6mm, bottom=1.6mm,
  fontupper=\ttfamily\small\color{qcode}}

\newtcolorbox{license}{skin=enhanced, colback=white, colframe=qgrey!55,
  boxrule=0.5pt, arc=0.6mm, left=3mm, right=3mm, top=2mm, bottom=2mm,
  fontupper=\small}

\setlist[itemize]{leftmargin=1.3em, itemsep=1.5pt, topsep=2pt}
\setlist[enumerate]{leftmargin=1.5em, itemsep=3pt, topsep=2pt}

\begin{document}

% ---------------- title page ----------------
\begin{titlepage}
\centering
\vspace*{3.2cm}
{\Huge\bfseries\color{qblue} qmeas}\\[6pt]
{\LARGE\color{qaccent} version 2}\\[10pt]
{\large user manual}\\[1.4cm]
{\normalsize\color{qgrey} Instrument control for cryogenic magneto-transport measurements}\\[2.5cm]
{\small\color{qgrey}
License and liability \textbullet{} installation \textbullet{} connecting instruments\\
the task grid \textbullet{} nesting and linking \textbullet{} while-loops \textbullet{} scripts \textbullet{} running}
\vfill
{\footnotesize\color{qgrey} Read the license and liability sections before relying on qmeas
for any measurement.\\ The rest is short --- skim \S3--\S4 once, then jump to whatever you need.}\\[1.4em]
{\footnotesize\itshape\color{qwarn} This manual was generated entirely by an AI system and may
contain errors.\\ Verify any critical detail against the software itself and your instruments'
own documentation.}\\[0.8em]
{\footnotesize\color{qgrey} qmeas and this manual were developed with Anthropic's Claude.}
\end{titlepage}

\tableofcontents
\newpage

% ======================================================================
\section{License}

qmeas is released under the \textbf{MIT License}. In plain terms: you may
use, copy, modify, and redistribute it freely, including in commercial and
closed-source work, provided the copyright notice and the license text below
are kept with any substantial copy. There is no warranty --- see the
liability section that follows, which is part of the license, not an
add-on.

\begin{license}
MIT License\\[4pt]

Copyright (c) 2026 Project h\_e2\\[4pt]

Permission is hereby granted, free of charge, to any person obtaining a copy
of this software and associated documentation files (the ``Software''), to
deal in the Software without restriction, including without limitation the
rights to use, copy, modify, merge, publish, distribute, sublicense, and/or
sell copies of the Software, and to permit persons to whom the Software is
furnished to do so, subject to the following conditions:\\[4pt]

The above copyright notice and this permission notice shall be included in
all copies or substantial portions of the Software.\\[4pt]

THE SOFTWARE IS PROVIDED ``AS IS'', WITHOUT WARRANTY OF ANY KIND, EXPRESS OR
IMPLIED, INCLUDING BUT NOT LIMITED TO THE WARRANTIES OF MERCHANTABILITY,
FITNESS FOR A PARTICULAR PURPOSE AND NONINFRINGEMENT. IN NO EVENT SHALL THE
AUTHORS OR COPYRIGHT HOLDERS BE LIABLE FOR ANY CLAIM, DAMAGES OR OTHER
LIABILITY, WHETHER IN AN ACTION OF CONTRACT, TORT OR OTHERWISE, ARISING
FROM, OUT OF OR IN CONNECTION WITH THE SOFTWARE OR THE USE OR OTHER DEALINGS
IN THE SOFTWARE.
\end{license}

% ======================================================================
\section{Liability and safety disclaimer}

The warranty disclaimer above is legal boilerplate; this section spells out
what it means for a program that drives real cryogenic and high-field
hardware.

\begin{warn}
qmeas sends commands to laboratory instruments --- magnets, sources,
cryostats --- that can damage equipment, samples, or people if operated
incorrectly. The software is provided without warranty of any kind. The
authors accept \textbf{no liability} for any loss, damage, injury, or cost
arising from its use, including from software defects, incorrect device
configuration, mis-entered command strings, network faults, or
instrument-specific behaviour the software cannot anticipate.
\end{warn}

Practical consequences, all of which are your responsibility, not the
software's:

\begin{itemize}
  \item \textbf{You are responsible for every command string you define.}
  qmeas transmits what you tell it to. A wrong SCPI string, a wrong ramp
  rate, or a wrong verify condition can drive an instrument outside safe
  limits. Verify command strings against the instrument's own manual.
  \item \textbf{Hardware interlocks are not qmeas's job.} Magnet quench
  protection, temperature limits, current limits, and emergency stops must
  be enforced by the instruments and their own safety systems. Do not rely
  on qmeas as a safety layer.
  \item \textbf{Test new sequences with Simulate, and at low/safe setpoints
  first.} Simulate shows what a task list will do without touching hardware.
  Use it before any long or high-stakes run. It is a convenience, not a
  guarantee.
  \item \textbf{Supervise runs that move real fields or temperatures.} The
  \code{onhold} and Stop mechanisms are best-effort; a network drop or a
  wedged instrument can prevent a clean abort.
\end{itemize}

By installing and using qmeas you accept these terms.

% ======================================================================
\section{Installation}

qmeas runs on \textbf{Windows} (the lab's target platform) and needs Python
plus a few packages. The optional AI Assistant additionally needs Ollama;
that part can be skipped entirely --- the measurement software is fully
functional without it.

\subsection{Step 1 --- Python}

Install \textbf{Python 3.11 or newer} (64-bit) from
\href{https://www.python.org/downloads/}{python.org}. The python.org
installer also installs the \textbf{Windows Python launcher} \code{py},
which every command in this manual uses --- it finds the right Python via
the registry, so it works even when \code{python} on PATH points at another
installation (Anaconda, an old version) or at nothing. (Ticking ``Add Python
to PATH'' during install is still fine, just no longer load-bearing.)

Confirm it in a \emph{new} PowerShell or Command Prompt window:

\begin{cmd}
py --version
\end{cmd}

\subsection{Step 2 --- Python packages}

qmeas depends on three third-party packages; everything else it uses is in
Python's standard library. Install them with pip \emph{through the launcher}:

\begin{cmd}
py -m pip install wxpython matplotlib pyvisa pyvisa-py
\end{cmd}

Always use \code{py -m pip}, not bare \code{pip}, and launch qmeas with
\code{py}: the two then use the \emph{same} Python. On machines with several
Pythons (Anaconda, old installs), bare \code{pip} and \code{py} can point at
different interpreters --- the classic symptom is
\code{ModuleNotFoundError: No module named 'wx'} at startup even though the
install ``worked''.

What each is for:

\begin{center}\small
\begin{tabular}{@{}p{2.9cm}p{10.7cm}@{}}
\toprule
\code{wxpython} & The GUI toolkit --- all windows, panels, the task grid.
  Required. \\
\code{matplotlib} & The live Graph window. If it's missing, qmeas still
  runs and measures; only the plot window shows a placeholder. \\
\code{pyvisa} & The instrument-communication layer (GPIB, USB, serial,
  TCPIP). Required for talking to any real device. \\
\code{pyvisa-py} & Pure-Python VISA backend, used automatically as a
  fallback when no vendor VISA library (NI-VISA) is present. Covers
  TCPIP/socket instruments --- which is most of this lab. \\
\bottomrule
\end{tabular}
\end{center}

\begin{warn}
\textbf{GPIB or USB instruments will not work until you additionally install
the NI-VISA library
(\href{https://www.ni.com/en/support/downloads/drivers/download.ni-visa.html}{download
from ni.com}).} The pure-Python \code{pyvisa-py} backend only handles
TCPIP/socket devices but not GPIB or USB! If qmeas can't load a vendor
library, it falls back to \code{pyvisa-py} automatically and warns you once.
\end{warn}

\subsection{Step 3 --- Run qmeas}

Put the qmeas files (\code{qmeas.py} and its companion modules, the
\code{images/}, \code{devices/}, and \code{custom/} folders) in one folder,
open a terminal there, and run:

\begin{cmd}
py qmeas.py
\end{cmd}

That's the whole measurement install. The next section is only needed for
the optional AI Assistant.

% ======================================================================
\section{Optional --- the AI Assistant (Ollama)}
\label{sec:assistant}

qmeas includes a built-in AI assistant, opened via \textbf{View $\rightarrow$
AI Assistant}. It can answer questions about qmeas --- what a red LED means,
how nesting works, why a row failed --- and help you set up an experiment:
which device settings, commands, and task rows a measurement needs. It is
optional; qmeas measures perfectly well without it.

The assistant runs a language model \textbf{locally} through Ollama --- nothing
is sent to any cloud service, which is the point: measurement data and device
configurations never leave the machine.

\begin{warn}
This is a \textbf{local model with limited capabilities} --- far smaller than
cloud AI services --- and it \textbf{may give incorrect answers}, including
confidently wrong ones. Treat its output as a starting point, not an
authority: check suggested device settings and command strings against the
instrument's own manual, and use Simulate before running anything it helped
set up.
\end{warn}

\subsection{Step 1 --- Install Ollama}

Download the Windows installer from
\href{https://ollama.com/download/windows}{ollama.com/download/windows} and
run it (no administrator rights needed for a standard install). It installs a
background service and a system-tray app, and exposes a local API at
\code{http://localhost:11434}. If Windows Firewall prompts on first start,
allow access for \textbf{private} networks --- this is local only, nothing is
exposed to the internet.

Confirm in a new terminal:

\begin{cmd}
ollama --version
\end{cmd}

\subsection{Step 2 --- Pull a model}

Models are downloaded with \code{ollama pull}. The recommended model is
\code{llama3.1:8b} --- it fits an 8~GB GPU and, in testing, follows the
assistant's structured instructions (e.g.\ the two ways to set up a magnet's
verify) noticeably more reliably than similarly-sized alternatives:

\begin{cmd}
ollama pull llama3.1:8b
\end{cmd}

Test that it responds:

\begin{cmd}
ollama run llama3.1:8b "Reply with the single word: ok"
\end{cmd}

If you have a larger GPU (or a remote GPU machine), a bigger model gives
better answers on multi-step setup questions --- see the table below. What
matters most for this assistant is instruction-following, not raw size, so a
strong 8B is a good default rather than a compromise.

\subsection{Step 3 --- Hardware and model size}

Local models are limited by GPU memory (VRAM). Rough guide for quantised
(Q4) models:

\begin{center}\small
\begin{tabular}{@{}p{3.6cm}p{4.6cm}p{5.4cm}@{}}
\toprule
\textbf{GPU VRAM} & \textbf{Comfortable model} & \textbf{Notes} \\
\midrule
8 GB  & 8B (\code{llama3.1:8b}) & Recommended; fits entirely in VRAM and
  answers fluently. \\
12 GB & 13--14B & A step up if it fits fully in VRAM. \\
24 GB & 32B & Good middle ground. \\
48 GB+ & 70--72B (e.g.\ \code{llama3.1:70b}) & Best answers on
  multi-step setup questions. \\
\bottomrule
\end{tabular}
\end{center}

Key rule: pick a model that fits ENTIRELY in your GPU's VRAM. A model that
spills into system RAM/CPU (e.g.\ a 14B on an 8~GB card, or a 70B without the
VRAM for it) becomes drastically slower --- usually worse than a smaller model
that fits. On the common 8~GB lab cards (e.g.\ NVIDIA T1000), \code{llama3.1:8b}
is the right choice; a 14B is not worth it there.

No suitable GPU? Models still run on CPU, just slowly --- fine for occasional
questions, tedious for anything interactive. A large model on CPU is
impractically slow, so use an 8B there.

\begin{tip}
Ollama unloads a model from memory after 5 minutes idle by default, so the
next question waits for a reload. If the assistant feels sluggish between
questions, set the environment variable \code{OLLAMA\_KEEP\_ALIVE} to
\code{30m} (or longer) to keep the model resident.
\end{tip}

\subsection{Step 4 --- Use it}

Open \textbf{View $\rightarrow$ AI Assistant}. qmeas connects to Ollama at
its default local address (\code{localhost} port \code{11434}) and lets you
choose among the models you've pulled. No further configuration is needed
beyond having Ollama running with at least one model downloaded.



% ======================================================================
\section{The five panels}

qmeas is a set of dockable panels. Drag any panel's title bar to move or
undock it; the \textbf{View} menu shows/hides panels and offers vertical or
horizontal stacking presets.

\begin{itemize}
  \item \textbf{Devices \& Commands} --- the tree of instruments and the
  read/write commands you've defined for each. Double-click a command to
  drop it into the task list.
  \item \textbf{Device Editor} --- add or edit an instrument's address and
  connection settings. Hidden until you click a device or open it from
  \textbf{View}.
  \item \textbf{Tasks} --- the ordered grid of steps that make up a run.
  Most of this manual is about building rows here.
  \item \textbf{Log} --- a timestamped record of each executed row.
  \item \textbf{Graph} --- a \emph{separate window} (not a docked panel, so
  you can throw it on a second monitor) that live-plots the running sweep.
\end{itemize}

Two virtual entries always sit in the device tree: \code{control} (pause,
prompt, stop, counter, and while-loops --- \S\ref{sec:while}) and \code{script} (runs a
Python file from the \code{custom/} folder; off by default).

% ======================================================================
\section{Connecting a device}

This is where most first-time trouble lives, so it gets the most space.
Open the \textbf{Device Editor}, click \textbf{New Device}, and work through
the buttons.

\subsection{Scan first --- but don't trust an empty result}

\textbf{Scan} asks the VISA layer to list everything it can see: GPIB, USB,
serial, and instruments that answer VISA's LAN discovery.

\begin{warn}
A device that does \emph{not} appear in the scan is very often still there
and perfectly usable. Most LAN instruments in this lab (lock-ins, magnet
controllers, source-measure units) speak over a plain \textbf{raw TCP
socket} and have no way to announce themselves --- the scan simply can't
see them. \textbf{This is normal.} You add them by hand (next section). An
empty or short scan list is not a fault and not a reason to stop.
\end{warn}

\subsection{Adding a socket (LAN) device by IP}\label{sec:connecting}

Click \textbf{Add TCPIP\ldots}, enter the instrument's IP address and port,
and choose \emph{Raw socket (SOCKET)} --- the common case here. (Choose
\emph{Standard VISA (INSTR)} only if the instrument's manual explicitly
documents VXI-11 or HiSLIP.)

You must know the IP before this step; qmeas cannot discover it for a
socket device. The networking pitfalls, in rough order of how often they
bite:

\begin{itemize}
  \item \textbf{Same subnet, or nothing connects.} The instrument and the
  qmeas PC must be on the same local network --- in practice the first three
  numbers of their IPs match, e.g.\ both \code{192.168.1.x}. If they differ
  (\code{192.168.1.x} vs \code{192.168.2.x}), the connection just times out.
  That's a network problem, not a qmeas problem, and no setting inside qmeas
  will fix it.
  \item \textbf{DHCP moves addresses.} Most instruments ship set to
  automatic (DHCP) addressing. Convenient --- until the instrument is
  power-cycled and comes back on a \emph{different} IP, silently breaking a
  device that worked yesterday. For anything permanent, either set a static
  IP in the instrument's own network menu, or reserve its address on the lab
  router (keyed to its MAC). Then it never moves again.
  \item \textbf{Find the current IP} from the instrument's front-panel
  network/LAN menu, or your router's client list.
  \item \textbf{Ping before qmeas.} From the qmeas PC, run \code{ping
  <ip>}. No reply $\Rightarrow$ fix the network first; qmeas will do no
  better than ping can.
\end{itemize}

\begin{gridex}{Example --- a magnet controller on a raw socket}
IP\ \ \ \ 192.168.1.42\\
Port\ \ 7020\\
Type\ \ Raw socket (SOCKET)
\end{gridex}

\subsection{HTTP / JSON-RPC devices}

Some instruments expose an HTTP server rather than a socket (e.g.\ the
Kiutra cryostat's JSON-RPC API). Click \textbf{Add HTTP\ldots} and give the
host, port, path, and \code{Content-Type}. Each command is then sent as an
HTTP POST with your command string as the body; the response body comes back
as the reading. The command strings themselves (the JSON-RPC or REST text)
are defined later in the Command Editor, exactly like SCPI strings.

\begin{gridex}{Example --- Kiutra JSON-RPC}
Host\ \ \ \ \ \ \ \ \ \ 192.168.1.55\\
Port\ \ \ \ \ \ \ \ \ \ 1006\\
Path\ \ \ \ \ \ \ \ \ \ /\\
Content-Type\ \ application/json-rpc
\end{gridex}

\subsection{Termination characters matter}

Every instrument expects each message to end with specific character(s) ---
the \textbf{write/read termination}. Get these wrong and the symptom is
maddening: the connection opens fine, but every command times out or returns
nothing, because the instrument is still waiting for an end-of-message it
never received. Most raw-socket SCPI gear here uses \code{\textbackslash n}
for both. When in doubt, the instrument's programming manual states it
explicitly --- check it before assuming the device is broken.

\subsection{When it still won't talk}

\begin{tip}
Before deeper debugging: \textbf{restart the instrument}. Single-client LAN
instruments can get stuck holding a half-open connection from a previous
crash or an ungraceful disconnect, and will refuse or drop new connections
until power-cycled. A restart clears this surprisingly often --- try it early,
not as a last resort.
\end{tip}

\subsection{Devices without a network interface (qbridge)}
\label{sec:qbridge}

Some instruments have \emph{no} instrument interface at all --- no SCPI, no
socket, no VISA. The Quantum Design OptiCool is the standing example: it is
controlled exclusively through its MultiVu software's OLE interface on the
control PC. qmeas talks to such devices through a \textbf{qbridge}: a small
translator process, shipped with qmeas in the \code{qbridge/} folder, that
speaks the vendor's interface on one side and plain socket text on the other.
To qmeas, the result is an ordinary raw-socket device.

\textbf{OptiCool setup} (once, on the OptiCool control PC):

\begin{enumerate}
  \item MultiVu running, as always.
  \item Install Quantum Design's Python package and start its server next to
  MultiVu: \code{py -m pip install MultiPyVu} then \code{py -m MultiPyVu}.
  \item Start the bridge: \code{py opticool.py} (from the \code{qbridge/}
  folder; listens on port 5100). Set both to auto-start with Windows so you
  never think about them again.
\end{enumerate}

Then, in qmeas on any lab PC: \textbf{Add TCPIP\ldots} with the control PC's
IP, port \code{5100}, raw socket, terminators \code{\textbackslash n} ---
exactly like any other LAN instrument (\S\ref{sec:connecting} applies,
including Test Connection, which the bridge answers). Define commands from
this verb table:

\begin{center}\small
\begin{tabular}{@{}p{5.2cm}p{8.6cm}@{}}
\toprule
\textbf{Command string} & \textbf{Meaning} \\
\midrule
\code{TEMP?} & temperature (K), e.g.\ \code{4.2130} \\
\code{TSTATE?} & temperature status: \code{Stable}, \code{Tracking},
  \code{Near}, \code{Chasing}, \code{Standby} \\
\code{AUXTEMP?} & auxiliary thermometer (K) \\
\code{FIELD?} & field (Oe; 10000\,Oe = 1\,T) \\
\code{FSTATE?} & field status: \code{Holding (driven)}, \code{Ramping}, \ldots \\
\code{TEMP [\%],10} & set temperature to the row value, 10\,K/min \\
\code{FIELD [\%],110} & set field to the row value (Oe), 110\,Oe/s \\
\code{FHOLD} & freeze the magnet at its current field --- use as this
  device's \code{onhold} command so Stop/abort parks the magnet \\
\bottomrule
\end{tabular}
\end{center}

Everything then works as for any device: sweep \code{FIELD [\%],110} rows,
log \code{TEMP?} green during any sweep, \emph{verify} a field step on
\code{FSTATE?} equals \code{Holding (driven)} (type the state string
exactly), or gate with a while-row on \code{TEMP?}. Waiting for stability is
deliberately done in qmeas --- visible, abortable, logged --- not inside the
bridge.

\begin{tip}
If OptiCool reads fail (\code{READ FAILED} in the data), the usual suspects
in order: the bridge process isn't running; the MultiPyVu server isn't
running; MultiVu isn't running. Restart in that order --- the bridge retries
its upstream connection once by itself, so a brief MultiPyVu hiccup usually
heals without intervention.
\end{tip}

\begin{note}
Two OptiCool adapter variants ship, serving identical commands:
\code{opticool.py} (via Quantum Design's MultiPyVu --- preferred, but needs
current OptiCool software) and \code{opticool\_qdi.py} (via the QDInstrument
.NET server on port 11000, the LabVIEW interface --- works on older software,
can run on the qmeas PC itself, needs \code{py -m pip install pythonnet} and
\code{QDInstrument.dll} beside it). Only the bridge changes; qmeas
configuration is the same for both. State strings differ between the two ---
always Query Once before using one as a verify target.
\end{note}

A third shipped adapter, \code{toptica\_bridge.py}, covers TOPTICA DLC~pro/DLC~smart
controllers (TeraScan): their command line is session-oriented (telnet banner
plus an access-level handshake), so the bridge holds one SDK session and
exposes the whole DeCoF parameter tree generically --- qmeas command strings
\code{GET <param>}, \code{SET <param> [\%]}, \code{EXEC <cmd>} reach any
parameter without adapter changes (\code{py -m pip install toptica\_lasersdk}; run
\code{py toptica\_bridge.py --toptica-host <ip>}). Verified TeraScan (DLC
smart THz) commands, units GHz:

\begin{center}\small
\begin{tabular}{@{}p{4.6cm}p{9.2cm}@{}}
\toprule
\code{GET frequency:frequency-act} & actual THz frequency --- log it green:
  the readback honestly wanders at the MHz level (the THz output is a
  two-DFB-laser beat), so recording the actual beats trusting the setpoint \\
\code{SET frequency:frequency-set [\%]} & frequency setpoint; use a FLOAT
  number format (DeCoF real-typed parameters reject bare integers) \\
\code{GET frequency:fast-scan-isscanning} & 1 while the native fast scan
  runs, else 0 \\
\bottomrule
\end{tabular}
\end{center}

Point-by-point stepping (a \code{SET \ldots frequency-set [\%]} sweep row
with the reads green) is the qmeas-idiomatic mode and composes with nesting;
the native continuous TeraScan sweep does not map onto task rows. Discover
further parameter names with the \code{decof\_tree\_dump.py} helper or the
TeraScan software's parameter reference. Note that each qmeas read opens a
short-lived connection \emph{to the bridge} (visible in its log) --- normal;
the bridge's own session to the controller persists.

Other interface-less instruments follow the same pattern: copy an adapter
file in \code{qbridge/}, map a handful of verbs to the vendor's API, and the
device becomes a socket instrument.

\subsection{Finishing up}

Fill in the \textbf{Alias} (lowercase letters/digits, e.g.\ \code{magnet};
every command becomes \code{alias\_\allowbreak commandname}), termination, timeout
(milliseconds; 5000 is fine for almost everything), and the device list it
joins. Click \textbf{Test Connection} --- for SCPI it sends \code{*IDN?} and
shows the reply; for HTTP it checks the host/port are reachable. Then
\textbf{Save Device}.

% ======================================================================
\section{Commands}

A new device has no commands. Right-click it $\rightarrow$ \textbf{Add
Command\ldots} A \textbf{write} sends a value; \code{[\%]} in its command
string is where qmeas substitutes the value from the task row. A
\textbf{query} reads a value back; use \textbf{Query Once} to see the raw
reply and click-select the part you want.

\begin{gridex}{Example --- two commands on a source aliased \texttt{gate1}}
Write:\ \ name gate1\_setvoltage\ \ \ \ string  VOLT [\%]\\
Query:\ \ name gate1\_readcurr\ \ \ \ \ \ string  CURR?
\end{gridex}

\subsection{Write methods: none, ramp, verify}\label{sec:methods}

A write command can carry a \textbf{Method}, chosen from a dropdown in the
Command Editor (visible only for a write command that contains a \code{[\%]}
placeholder). It governs \emph{how} the value is delivered.

\textbf{none} --- write the value in one shot and move on. The default;
correct for almost everything.

\textbf{ramp} --- approach the target gradually instead of jumping. qmeas
itself steps the value in one-second increments, each no larger than the
rate you set, all the way from the current value to the target. Use it
whenever a sudden jump would be unsafe or unphysical (a gate that must not
slew, say). Two fields appear:

\begin{itemize}
  \item \textbf{Ramp rate (units/s)} --- the maximum change per second, in
  the same units as the command's value. Must be greater than zero.
  \item \textbf{Based on output state of} --- a read-back (query) command on
  the \emph{same} device, used to find the \emph{starting} value so qmeas
  knows where the ramp begins. Pick the query that reports this instrument's
  present output.
\end{itemize}

\begin{warn}
Ramp is \textbf{write-and-forget}: qmeas issues the intermediate setpoints on
a one-second clock but does \emph{not} check the instrument actually reached
each one. It paces commands; it does not confirm the hardware kept up. If the
instrument has its own internal ramp/slew feature, decide which one owns the
rate --- don't set both against each other. The ramp rate is your
responsibility to keep within the instrument's safe limits.
\end{warn}

\textbf{verify} --- after writing, repeatedly read a chosen read-back until it
reports the expected value, or give up. This is how a step waits for hardware:
a field setpoint that blocks until the magnet controller reports \code{HOLD},
for instance. Fields:

\begin{itemize}
  \item \textbf{Based on output state of} --- the read-back (query) command to
  poll, on the same device.
  \item \textbf{\ldots equals} --- the value that read-back must report for the
  step to be considered done. For a status string, type it exactly
  (\code{HOLD}); for a number, the target value.
  \item \textbf{Tolerance ($\pm$, 0 = exact)} --- how close a numeric read-back
  must be to count as a match. Leave at \code{0} for an exact match (the right
  choice for status strings like \code{HOLD}); widen it for a numeric setpoint
  that will never land on the exact digits (e.g.\ a field good to $\pm$0.0005\,T).
  \item \textbf{Verify timeout (s, 0 = none)} --- how long to keep polling
  before giving up. On timeout the run \textbf{aborts} (with magnet onhold);
  \code{0} means wait indefinitely.
\end{itemize}

\begin{gridex}{Example --- a field setpoint that waits for HOLD}
Command:\ \ magnet\_setfield\ \ \ \ \ \ string  SOUR:FIELD [\%]\\
Method:\ \ \ verify\\
Based on output state of:\ \ magnet\_status\\
...equals:\ \ HOLD\ \ \ \ \ \ Tolerance:\ 0\ \ \ \ \ Verify timeout:\ 600
\end{gridex}

\noindent Put that command in a task row with the field value in
Constant/Start, and the row won't finish --- and the next row won't begin ---
until \code{magnet\_status} reads \code{HOLD} or 600\,s elapse.

\begin{tip}
\textbf{Ramp vs.\ verify are independent, not two halves of one thing.} Ramp
paces the \emph{outgoing} setpoints; verify waits on the \emph{incoming}
read-back. A magnet controller with its own internal ramp is the classic
verify case: you write the setpoint, let the controller ramp on its own, and
\emph{verify} that it reports \code{HOLD} when done --- no qmeas ramp involved.
A source with no internal ramp is the classic ramp case. You pick one method
per command.
\end{tip}

\subsection{Query LEDs: green, red, grey}

Every query command shows a small coloured LED next to it in the Devices \&
Commands tree. Click it to cycle through three states, which control whether
that reading is taken during a run and whether it's saved:

\begin{itemize}
  \item \textbf{Green --- read and save.} The query is polled at every
  measurement step and its value is written to the data file. This is the
  normal state for any reading you want to record (a lock-in output, a field
  value).
  \item \textbf{Red --- read for verification, don't save.} The query is still
  available for a write command's \emph{verify} method (\S\ref{sec:methods}) to
  poll, but it is \emph{not} added to the data file as a measured column. Use
  it for a status read-back --- e.g.\ \code{magnet\_status} watched by a verify
  step --- that you don't want cluttering the saved data with a \code{HOLD}
  column.
  \item \textbf{Grey --- off.} The query is not polled at all and not saved.
\end{itemize}

Red and grey behave identically as far as the data file is concerned (neither
is saved); the difference is a reminder to you that a red query is deliberately
in use \emph{somewhere} as a verify target, whereas grey is simply idle. The
state is remembered when you close qmeas and restored next time.

\begin{tip}
If a status string like \code{HOLD} or a \code{READ FAILED} message is showing
up as a column in your data files and you don't want it there, the query is
green --- click its LED to red. It stays available to any verify step that uses
it, but drops out of the saved data.
\end{tip}

% ======================================================================
\section{The task grid}

Each row is one step, run top to bottom. Columns:

\begin{center}\small
\begin{tabular}{@{}p{2.6cm}p{11.3cm}@{}}
\toprule
\textbf{On} & Checkbox. Unchecked rows are skipped, not deleted. \\
\textbf{Command} & \code{device\_command}, dragged from the tree. \\
\textbf{Constant/Start} & The single value (constant write) or the sweep's
  first value. On a linked row, a \code{[\%]} expression instead (\S\ref{sec:linking}). \\
\textbf{Final} & The sweep's last value. Blank = constant write. \\
\textbf{Steps} & Number of points, Start$\to$Final. Blank/\code{0} =
  constant. \\
\textbf{Integration Time} & Wait after each step: \code{2s}, \code{0.5s},
  \code{2m} (bare number = seconds; \code{ms} not accepted --- use \code{0.5}
  not \code{500ms}). \\
\textbf{Comment} & Free text, copied into the Log. \\
\bottomrule
\end{tabular}
\end{center}

\begin{tip}
Fields turn \textcolor{qwarn}{\textbf{red}} when a row is inconsistent ---
e.g.\ you filled Final but left Steps or Integration Time empty, or typed a
non-number. A red cell never runs the way you'd want; fix it before pressing
Start. (This is why the earlier \code{10a}-in-Steps class of silent mistake
can't slip through unseen.)
\end{tip}

The \code{control} device gives four fixed rows: \code{pause} (waits the
Constant value in seconds), \code{userprompt} (popup, waits for you),
\code{stop} (aborts, same as the Stop button), and \code{counter} (a
step/time index that writes no instrument --- behaves like a sweep for
nesting and linking).

% ======================================================================
\section{Nesting: sweeps inside sweeps}

Nesting chains rows into nested for-loops: for each step of the outer row,
the inner row runs its whole sweep. Select the row that should be inner,
right-click $\rightarrow$ \textbf{Nest}; it attaches to the row directly
above, one level deeper. Build deeper chains bottom-up, one Nest at a time.
Max depth 4. \textbf{Unnest} peels from the bottom.

\begin{gridex}{Example --- gate sweep at every field}
Row 1\ \ magnet\_setfield\ \ \ \ Start 0\ \ Final 1\ \ Steps 6\ \ \ IntTime 5s\\
Row 2\ \ \ \ gate1\_setvoltage\ Start -2\ Final 2\ Steps 21\ IntTime 0.2s\ \ (nested)
\end{gridex}

\noindent Right-click Row 2 $\rightarrow$ Nest. The run sets field to 0\,T,
sweeps the gate through all 21 points, steps the field to 0.2\,T, sweeps the
gate again, \ldots\ $6\times21=126$ points. A counter can be the outer loop
too --- ``every 10\,s, sweep the gate'' for a drift study:

\begin{gridex}{Example --- counter as the outer (time) loop}
Row 1\ \ control\_counter\ \ \ \ \ Steps 100\ \ IntTime 10s\\
Row 2\ \ \ \ gate1\_setvoltage\ Start -2\ Final 2\ Steps 21\ IntTime 0.2s\ \ (nested)
\end{gridex}

% ======================================================================
\section{Linking: rows that track a driver}\label{sec:linking}

Linking makes ``follower'' rows compute their value from a ``mother'' row's
swept value, instead of sweeping on their own. Select the follower,
right-click $\rightarrow$ \textbf{Link}; it attaches to the row above. Its
Constant/Start holds an \textbf{expression} with \code{[\%]} standing for the
mother's current value. Operators \code{+ - * / \textasciicircum} and
functions \code{sin cos tan exp log log10 sqrt abs} (and \code{pi},
\code{e}) are available.

\begin{gridex}{Example --- symmetric split gates (gate2 = -gate1)}
Row 1\ \ gate1\_setvoltage\ \ Start -2\ Final 2\ Steps 41\ IntTime 0.3s\ \ (mother)\\
Row 2\ \ \ \ gate2\_setvoltage\ Start [\%]*-1\ \ \ \ \ \ \ \ \ \ \ \ \ \ \ \ \ \ \ \ \ (follower)
\end{gridex}

\noindent At the step where \code{gate1}\,=\,+0.5\,V, \code{gate2} is written
to $-0.5$\,V before the reading. A follower can itself use a ramp/verify
method and a nonlinear expression:

\begin{gridex}{Example --- heater tracking a time index nonlinearly}
Row 1\ \ control\_counter\ \ \ \ \ Steps 50\ IntTime 1s\ \ (mother)\\
Row 2\ \ \ \ heater\_setpower\ \ Start 0.5*sin([\%]/8)\textasciicircum 2\ \ (follower)
\end{gridex}

\begin{tip}
A follower's Final, Steps, and Integration Time are driven by the mother, so
they're locked. Its Constant/Start (the expression) is \emph{required} --- an
empty one turns red.
\end{tip}

% ======================================================================
\section{What combines with what}

Mother rows and nesting parents must have a value to sweep; the
value-less row types (while, script, pause/prompt/stop) can only sit as
plain standalone rows. Nesting and linking never combine on the same row.
Everything below is enforced both in the menus (grayed out) and at run time.

\begin{center}\small
\begin{tabular}{@{}p{4.5cm}*{5}{>{\centering\arraybackslash}p{1.55cm}}@{}}
\toprule
\textbf{Row type} & \textbf{Nest parent} & \textbf{Nest child} &
  \textbf{Max depth} & \textbf{Link mother} & \textbf{Link follower} \\
\midrule
Device write (Steps $\geq$ 1) & \checkmark & \checkmark & 4 & \checkmark & \checkmark \\
\code{control\_counter}       & \checkmark & \checkmark & 4 & \checkmark & $\times$ \\
Virtual while-loop            & $\times$   & $\times$   & --- & $\times$ & $\times$ \\
\code{script\_*}              & $\times$   & $\times$   & --- & $\times$ & $\times$ \\
\code{pause}/\code{userprompt}/\code{stop} & $\times$ & $\times$ & --- & $\times$ & $\times$ \\
\bottomrule
\end{tabular}
\end{center}

\noindent The value-less rows compose by \emph{sequence}, not structure ---
just put them one after another as ordinary rows:

\begin{gridex}{Example --- gate a nested sweep behind a hardware wait}
Row 1\ \ magnet\_setfield\ \ \ \ \ Start 1\ \ (method verify: waits for readback = 1 T)\\
Row 2\ \ control\_waitforhold\ \ (while: magnet\_status = HOLD, timeout 600s)\\
Row 3\ \ control\_counter\ \ \ \ \ Steps 100\ IntTime 10s\\
Row 4\ \ \ \ gate1\_setvoltage\ Start -2\ Final 2\ Steps 21\ IntTime 0.2s\ \ (nested)\\
Row 5\ \ script\_saveplot
\end{gridex}

% ======================================================================
\section{Virtual while-loops}\label{sec:while}

A while-loop is a \code{control} command that repeats until a condition on
some reading is met, then continues --- like verify, but as its own row that
also \emph{records data every poll}. Right-click \code{control}
$\rightarrow$ \textbf{Add Virtual While\ldots}, pick a query to watch (from
any device), an operator (\code{= != < >}), a value, and a timeout
(mandatory; \code{0} = wait forever). Drop it in the grid; its Integration
Time becomes the poll interval (minimum 0.1\,s).

\begin{gridex}{Example --- wait for base temperature, logging as it goes}
Definition:\ watch cryostat\_temperature, operator <, value 0.05, timeout 0\\
Grid row:\ \ \ control\_waitforbase\ \ \ IntTime 5s
\end{gridex}

\noindent Every 5\,s it reads the temperature (and all other active
queries, saved to the data file) and exits once it drops below 50\,mK.

\begin{warn}
Any while-loop failure --- timeout reached, a non-numeric reading against
\code{<}/\code{>}, the watched device switched off --- \textbf{aborts the whole
run} (with magnet onhold), by design: the row exists to gate progress on real
hardware state, so if it fails, later rows would run under unknown conditions.
\end{warn}

% ======================================================================
\section{Running, simulating, saving}

\textbf{Start} runs the list; \textbf{Stop} aborts (holding magnet devices).
\textbf{Simulate} walks the list and prints, in plain language, what each row
would do --- sweep ranges, nesting/linking with computed follower values,
while conditions --- \emph{without touching hardware}. It applies the same
rules as a real run, so it also flags invalid combinations. Use it to sanity
-check a list before a long run; it's optional, not required. \textbf{Save} /
\textbf{Load} store and restore a task list. Each sweep writes its own data
file, which the Graph window plots live.




\end{document}
